\documentclass[11pt]{article}
\usepackage[T1]{fontenc}
\usepackage[textwidth=250pt]{geometry}
\usepackage{amsmath,amssymb}
\usepackage{microtype}
\pagestyle{empty}
\begin{document}
The squares satisfy $x^2+y^2=z^2$ for the triple $(3,4,5)$, and the sum
$\sum_{i=1}^n i=\frac{n(n+1)}{2}$ follows by induction on $n\ge 1$. The
fraction $\frac12$ and the root $\sqrt{2}$ are irrational only in the second
case, and $a_1+a_2+\cdots+a_n$ is written $\sum a_i$ for short.

For every $n\ge 2$ the number $2^n-1$ is odd, and $3^2=9$ while $2^3=8$, so
the powers $x_1^2+x_2^2$ grow quickly with the index $i$ of each term; the
quotient $\frac{a+1}{b+2}$ is never an integer when $b=a+1$ holds.

The subscripts in $a_{n+1}=a_n+a_{n-1}$ define the sequence, and the ratio
$a_{n+1}/a_n$ tends to $(1+\sqrt5)/2$ as $n$ grows, which is the golden ratio
that appears in many problems of this particular kind in the course.
\end{document}
